\section{Theory}
\label{theory}

% ─────────────────────────────────────────────────────────────────────────
% PLAIN-LANGUAGE GLOSSARY for the technical terms in this section (comment
% only, not printed). Talking points for the supervisor meeting.
%   Crosslink: UV light "glues" the protein to the exact RNA spot it touches,
%       so sequencing can later read out that spot.
%   SMInput / input control: a background sample without the protein pulldown;
%       used to subtract non-specific signal (noise).
%   HMM (hidden Markov model): a statistical model that walks along the genome
%       and labels each position "bound" or "background"; PureCLIP uses one.
%   Bandwidth: how much PureCLIP smooths the signal. Merge distance: how close
%       two signals must be to count as one site. Both are tunable parameters.
%   PWM (position weight matrix): a table giving, per position, how much each
%       letter (A/C/G/U) is preferred. Used to score how well a site matches
%       the protein's motif. Log-odds = motif score vs. random-sequence score.
%   Shuffled background: randomly reordered sequence used as a "no real motif"
%       baseline, so we can tell true enrichment from chance.
%   Optuna / TPE (Tree-structured Parzen Estimator): a standard automated
%       search that tries parameter sets, learns from the scores, and proposes
%       the next set. No biology knowledge, purely data-driven. Our "baseline"
%       optimizer that the LLM is compared against.
% ─────────────────────────────────────────────────────────────────────────

\scaffold{Bullets to expand. Biology first, then computational tools; keep implementation in \S\ref{method}.}
\begin{itemize}
  \item eCLIP mechanism \parencite{vannostrand2016}: UV crosslink $\rightarrow$ immunoprecipitation $\rightarrow$ RT truncates at the crosslink $\rightarrow$ 5$'$ read end marks the contact at $\sim$nucleotide resolution. Two IP replicates + SMInput control.
  \item PureCLIP \parencite{krakau2017}: HMM over read 5$'$-start density; ``crosslinked'' vs.\ ``background'' states; optional SMInput covariate. Bandwidth = signal smoothing, merge distance = site grouping, wrong values give too few or too many sites.
  \item \textbf{Motif difficulty (background for one focused experiment).} Proteins differ in how sequence-defined their binding is: crisp k-mer motifs (RBFOX2 \texttt{UGCAUG}, PUM \texttt{UGUANAUA}) vs.\ degenerate (HNRNPK C-rich, SRSF1 GA-rich) vs.\ positional / non-k-mer (U2AF2 polypyrimidine tract, SF3B1/4 branch point). This spectrum is ``easy'' to ``hard'' for a motif-aware objective, and motivates a small side-experiment (\S\ref{results}) asking whether the LLM's prior helps more when the motif signal is weak. Note: the difficulty label is our manual annotation, not a quantity the pipeline computes.
  \item What makes a site set ``good'': reproducibility (chance-corrected), motif enrichment (PWM log-odds vs.\ shuffled background), benchmark recall (ENCODE reference regions).
  \item PWMs / position weight matrices: log-odds scoring; motifs drawn from several databases (mCrossBase, ATtRACT, RBPmap, oRNAment, RBPDB); the most-enriched motif is reported per site set.
  \item LLM as a reasoning optimizer (domain priors + reading structured feedback) vs.\ Bayesian optimization / TPE \parencite{bergstra2011}, as implemented in Optuna \parencite{akiba2019} (surrogate model, exploration/exploitation, no domain knowledge).
  % Figure to add (not printed): pipeline schematic eCLIP reads -> PureCLIP ->
  % sites -> scoring, plus a small motif-class table (crisp / degenerate / positional).
\end{itemize}

% Goal: give the reader the conceptual background needed to understand
% the project. Biological concepts first, then the computational tools.
% Implementation specifics belong in §3.

% ── 1. eCLIP and RNA-binding proteins ─────────────────────────────────
%   • RBPs regulate RNA post-transcriptionally (splicing, stability,
%     localisation, translation). Dysregulation is linked to cancer and
%     neurological disease.
%   • eCLIP (Van Nostrand et al. 2016): UV crosslink → immunoprecipitation
%     → sequencing. Reverse transcriptase truncates at the crosslink site,
%     so the 5′ read end marks the contact position at ~nucleotide
%     resolution. Two IP replicates + SMInput (size-matched input) control.
%   • The three RBPs studied, their canonical RNA motifs, and the two
%     cell lines (K562 leukemia, HepG2 liver). Brief table or inline list.

% ── 2. Peak calling with PureCLIP ─────────────────────────────────────
%   • Conceptual role: converts a pile-up of aligned reads into a discrete
%     set of binding sites. Necessary step between raw sequencing and
%     biological interpretation.
%   • PureCLIP mechanism: HMM over read 5′-start density; states
%     "crosslinked" vs. "background"; optionally uses SMInput as a
%     covariate to subtract non-specific signal.
%   • Why parameters matter: bandwidth controls smoothing of the signal;
%     merge distance controls how nearby sites are grouped. Wrong values
%     produce either too few (over-filtered) or too many (noisy) sites,
%     and the right values depend on the dataset.

% ── 3. What makes a binding-site set "good" ───────────────────────────
%   • Reproducibility: real binding recurs across independent replicates;
%     noise does not. Chance-correction is needed because broad intervals
%     overlap by chance even for random site sets.
%   • Motif enrichment: genuine sites should be enriched for the RBP's
%     known RNA recognition sequence relative to shuffled background.
%     PWMs (position weight matrices) provide a quantitative scoring model.
%   • Benchmark recall: good calls should recover the curated set of
%     strong regions from the ENCODE reference pipeline.

% ── 4. LLMs as reasoning-based optimizers ─────────────────────────────
%   • Large language models (LLMs) can incorporate domain knowledge
%     (priors about biology, known failure modes) into their proposals,
%     going beyond pure black-box function optimization.
%   • In an iterative loop, an LLM can read structured feedback (score
%     deltas per iteration) and reason about which parameter to change
%     next, analogous to a human expert tuning settings by hand, but
%     automated and consistent.

% ── 5. Bayesian optimization (Optuna / TPE) ───────────────────────────
%   • Bayesian optimization builds a probabilistic surrogate model of the
%     objective surface from previous evaluations and uses it to choose
%     the next point to evaluate (exploration vs. exploitation trade-off).
%   • Tree-structured Parzen Estimator (TPE) models the distribution of
%     good and bad parameter sets separately and samples from the ratio.
%     No domain knowledge needed; purely data-driven.
